\documentclass[11pt]{article}

%% Standard packages
\usepackage[letterpaper,left=1in,right=1in,top=1in,bottom=1.1in,includefoot,footskip=28pt]{geometry}
\usepackage{amsmath, amssymb}
\usepackage{mathtools}
\usepackage{bm}
\usepackage{bbm}
\usepackage{graphicx}
\usepackage{booktabs}
\usepackage{hyperref}
\usepackage{cleveref}
\usepackage{algorithm}
\usepackage{algorithmic}
\usepackage{xcolor}
\usepackage{microtype}
\usepackage{natbib}
\usepackage{fancyhdr}
\usepackage{titling}
\providecommand{\theHalgorithm}{\arabic{algorithm}}
\providecommand{\theHfigure}{\arabic{figure}}

%% Notation macros
\newcommand{\R}{\mathbb{R}}
\newcommand{\N}{\mathbb{N}}
\newcommand{\bx}{\mathbf{x}}
\newcommand{\bX}{\mathbf{X}}
\newcommand{\bz}{\mathbf{z}}
\newcommand{\bZ}{\mathbf{Z}}
\newcommand{\bd}{\mathbf{d}}
\newcommand{\bD}{\mathbf{D}}
\newcommand{\bB}{\mathbf{B}}
\newcommand{\bI}{\mathbf{I}}
\newcommand{\bJ}{\mathbf{J}}
\newcommand{\bU}{\mathbf{U}}
\newcommand{\bV}{\mathbf{V}}
\newcommand{\bW}{\mathbf{W}}
\newcommand{\norm}[1]{\left\|#1\right\|}
\newcommand{\abs}[1]{\left|#1\right|}
\newcommand{\innerprod}[2]{\langle #1, #2 \rangle}
\newcommand{\dgeo}{d^{\mathrm{geo}}}
\newcommand{\demb}{d^{\mathrm{emb}}}
\newcommand{\dknn}{d^{k\text{-NN}}}
\newcommand{\Gk}{G_k}
\newcommand{\Stress}{\mathrm{Stress}}
\newcommand{\StressOne}{S_1}
\newcommand{\sgn}{\operatorname{sgn}}
\newcommand{\diag}{\operatorname{diag}}
\newcommand{\placeholderfigurebox}[2][0.9\linewidth]{%
  \fbox{\parbox[c][0.18\textheight][c]{#1}{\centering\footnotesize #2}}%
}
\DeclareMathOperator*{\argmin}{arg\,min}
\setlength{\droptitle}{-2.2em}
\setlength{\headheight}{14pt}
\pagestyle{plain}

\newenvironment{keywords}{%
  \par\smallskip\noindent\textbf{Keywords: }\ignorespaces
}{\par\smallskip}

\newcommand{\paperaffiliationfooter}{%
  \parbox{0.92\textwidth}{\centering\footnotesize
  Center for Advanced Microbiome Research and Innovation (CAMRI), Institute for Genome Sciences, Department of Microbiology and Immunology, University of Maryland School of Medicine, Baltimore, Maryland, USA.}%
}

\fancypagestyle{firstpage}{%
  \fancyhf{}%
  \renewcommand{\headrulewidth}{0pt}%
  \renewcommand{\footrulewidth}{0pt}%
  \fancyfoot[C]{\paperaffiliationfooter}%
}

\providecommand{\manuscriptversion}{unversioned}
\providecommand{\manuscriptbuildnumber}{NA}
\providecommand{\manuscriptcommit}{unknown}
\providecommand{\manuscriptbuilddatetime}{unknown}
\InputIfFileExists{../build/manuscript_build_info.tex}{}{}

\begin{document}

\title{Geodesic Multidimensional Scaling:\\
       Self-Consistent Dimensionality Reduction\\ via Graph Geodesic Distances}

\author{Pawel Gajer \and Jacques Ravel}
\date{}

\maketitle
\thispagestyle{firstpage}
\begin{center}
\footnotesize
Version \texttt{\manuscriptversion} \quad | \quad Build
\texttt{\manuscriptbuildnumber} \quad | \quad Commit
\texttt{\manuscriptcommit}\\
Built \manuscriptbuilddatetime
\end{center}
\vspace{0.6em}

%% ============================================================
\begin{abstract}
Multidimensional scaling (MDS) seeks a low-dimensional embedding of
a set of points that faithfully reproduces a given pairwise dissimilarity
matrix.  Classical and metric MDS measure faithfulness via the discrepancy
between the input dissimilarities and the \emph{Euclidean} distances among
the embedded points.  Isomap improved upon MDS for structured data whose
relevant geometry is not well captured by ambient Euclidean distance by
replacing Euclidean input dissimilarities with shortest-path distances on a
$k$-nearest-neighbor ($k$-NN) graph, thereby approximating intrinsic
shortest-path geometry.  However, Isomap retains Euclidean distances on the
\emph{output} side, creating a fundamental asymmetry: input distances are
geodesic, but the fidelity criterion is Euclidean.  We argue that this
asymmetry is not merely an aesthetic inconsistency but a systematic source
of distortion whenever the data carry genuine graph, branching, or
bottleneck structure.

We propose \emph{Geodesic MDS}, a self-consistent formulation for
shortest-path metric spaces equipped with a chosen geodesic family.  The
input dissimilarities are geodesic distances, and the output fidelity criterion
measures the embedded lengths of the corresponding input-side geodesics.  The
embedding is sought in $\R^d$ such that a fixed family of shortest paths
chosen from the input $k$-NN graph retains its path-length geometry as
closely as possible after embedding.  We derive a closed-form
gradient for this geodesic stress function, extend the SMACOF majorization
framework to the geodesic setting, and characterize the role of $k$ as a
geometric scale parameter that continuously interpolates between a sparse
tree-like regime and standard metric MDS.  The resulting method is
practically implemented by replacing one line in a standard MDS pipeline
and adds no asymptotic complexity.  We situate the proposal within the
metric geometry literature by connecting it to minimum metric realizations,
tight spans, and metric spanners, and discuss applications to
high-dimensional compositional data, phylogenetic embeddings, and network
visualization.
\end{abstract}

\begin{keywords}
  multidimensional scaling, dimensionality reduction, geodesic distance,
  graph metric, $k$-nearest neighbor graph, Isomap, SMACOF, metric geometry,
  structured data geometry, stress function
\end{keywords}

%% ============================================================
\section{Introduction}
\label{sec:introduction}

The central task of \emph{multidimensional scaling} (MDS) is to find a
configuration of $n$ points $\bZ = (\bz_1, \ldots, \bz_n) \in \R^{n \times
d}$ in a low-dimensional Euclidean space $\R^d$ such that the Euclidean
distances $\norm{\bz_i - \bz_j}$ are as close as possible to a given $n
\times n$ symmetric, non-negative dissimilarity matrix $\Delta = (\delta_{ij})$
\citep{torgerson1952multidimensional, kruskal1964multidimensional,
borg2005modern, cox2000multidimensional}.  The criterion of closeness is made
precise by a \emph{stress} function, and the embedding is obtained by
minimizing this stress.  MDS and its relatives are among the oldest and most
widely used tools in multivariate statistics, visualization, and machine
learning \citep{shepard1962analysis, sammon1969nonlinear,
vandermaaten2008tsne, mcinnes2018umap}.

When the input data are given as a matrix $\bX \in \R^{n \times m}$ of $n$
observations in a high-dimensional ambient space $\R^m$, the most common
choice is $\delta_{ij} = \norm{\bx_i - \bx_j}_{\R^m}$, the Euclidean
distance in the original space.  This is appropriate only when ambient
Euclidean distance already reflects the meaningful geometry of the data.
But the fact that the observations are represented in $\R^m$ does not by
itself justify that choice: many real datasets have curved, branching,
bottlenecked, hole-rich, or otherwise strongly structured geometry for
which the relevant relationships are governed by paths through the data
rather than by straight lines in the ambient coordinates.  In such cases,
straight-line Euclidean
distances are poor proxies for the intrinsic distances that govern the
geometry of the data: two points near opposite ends of a curved surface
may be close in ambient Euclidean distance but far apart along the surface,
and the appropriate geodesic structure must be estimated from the data.

This limitation motivated \emph{Isomap}
\citep{tenenbaum2000global}, one of the most influential nonlinear
algorithms.  Isomap approximates the intrinsic geodesic distance between
points by constructing a $k$-nearest-neighbor ($k$-NN) graph on the data,
computing all-pairs shortest paths on this graph (weighted by Euclidean
edge lengths), and then feeding these graph-geodesic dissimilarities
$\delta_{ij} = \dknn(\bx_i, \bx_j)$ into classical MDS.  For data lying on
the special case of a smooth, sufficiently sampled Riemannian manifold, the graph-geodesic
distances converge to the true manifold geodesics as $n \to \infty$
\citep{bernstein2000graph, tenenbaum2000global}, and the resulting embedding
recovers the intrinsic geometry with high fidelity.

Despite this advance, Isomap retains a fundamental asymmetry that has not
previously been identified as a source of systematic distortion.  On the
\emph{input} side, distances are geodesic, meaning shortest paths through the data
graph.  On the \emph{output} side, the stress function measures faithfulness
via \emph{Euclidean} distances between embedded points.  This means that
Isomap optimizes an embedding in which the distance between $\bz_i$ and
$\bz_j$ is the straight-line distance $\norm{\bz_i - \bz_j}$, not the
embedded length of the corresponding input-side geodesic.

The comparison below summarizes the distinction between standard MDS,
Isomap, and the proposed Geodesic MDS in terms of the distance used on
the input and output sides of the embedding problem.

\begin{center}
  \small
  \begin{tabular}{@{}p{0.19\textwidth}p{0.34\textwidth}p{0.34\textwidth}@{}}
    \toprule
    Method & Input dissimilarity $\delta_{ij}$ & Output distance $d_{ij}$ \\
    \midrule
    Standard MDS & Euclidean $\norm{\cdot}_{\R^m}$ & Euclidean $\norm{\cdot}_{\R^d}$ \\
    Isomap & $d_{G_k}$ (graph geodesic) & Euclidean $\norm{\cdot}_{\R^d}$ \\
    Geodesic MDS & $d_{G_k}$ (graph geodesic) & embedded length of a fixed input geodesic \\
    \bottomrule
  \end{tabular}
\end{center}

For data whose intrinsic geometry is smooth and slowly varying, this
asymmetry causes only a small error: nearby embedded points have
approximately equal Euclidean and geodesic distances.  But for data with
genuine graph or tree structure, where the topology enforces that all
paths between certain pairs of points must pass through bottleneck
vertices, the Euclidean output distance is a systematic \emph{underestimate}
of the embedded length of the relevant graph geodesic (by the triangle inequality, with
equality only when points are collinear in the embedding space).

Consider a concrete example: a star graph with center $c$ and three arm
tips $v_1, v_2, v_3$, each arm of unit length.  The graph distance between
$v_1$ and $v_2$ is $\delta(v_1, v_2) = 2$ (through $c$).  Any faithful
embedding must respect the topological constraint that the only path from
$v_1$ to $v_2$ passes through $c$.  Standard MDS and Isomap both try to
place $\bz_{v_1}$ and $\bz_{v_2}$ at Euclidean distance $2$ apart, which
is perfectly achievable in $\R^2$, but doing so forces the embedded arms
to be nearly collinear, misrepresenting the star topology.  The \emph{correct}
criterion is that the \emph{path-length distance along the embedded arms}
from $\bz_{v_1}$ to $\bz_{v_2}$ equals $\norm{\bz_{v_1} - \bz_c} +
\norm{\bz_c - \bz_{v_2}} = 2$, which is compatible with any angular
arrangement of the arms and correctly reflects the star structure.
Figure~\ref{fig:star_bias} is intended to make this contrast visually explicit.

\begin{figure}[t]
  \centering
  \includegraphics[width=0.98\linewidth]{../figures/fig_star_bias.pdf}
  \caption{Branching metrics expose the output-side inconsistency.  In a star graph, the graph distance between leaves is realized through the center.  Matching this distance with Euclidean output distances biases the embedding toward collinearity, whereas measuring output fidelity by the embedded lengths of the corresponding input geodesics preserves the correct topological structure.}
  \label{fig:star_bias}
\end{figure}

This paper introduces \emph{Geodesic MDS}, a reformulation of the metric
MDS stress function that is self-consistent with respect to geodesic
distance.  We identify and formalize the input/output asymmetry in Isomap
and show that replacing Euclidean output distances with embedded lengths
of the input geodesics yields a self-consistent stress function
(\cref{sec:geodesic_mds}).  We define the geodesic stress function over a
fixed graph topology and a fixed family of input-side geodesics, where the
free variables are the vertex positions in the target space $\R^d$ and the
distances are measured along those embedded paths
(\cref{sec:geodesic_stress}).

We derive the gradient of the geodesic stress with respect to vertex
positions and characterize its non-local dependence on the chosen
input-side geodesics (\cref{sec:gradient}).  We extend the SMACOF
majorization framework
\citep{deleeuw1977applications, deleeuw2009multidimensional} to the
geodesic setting by exploiting the fact that the fixed input geodesics
decompose into sums of Euclidean edge lengths, enabling a tractable
iterative optimization via SMACOF.  We also characterize the
role of $k$ as a geometric scale parameter: $k=1$ (spanning tree)
corresponds to a tree-metric embedding, while $k \to n$ (complete graph)
recovers standard metric MDS (\cref{sec:properties}).  Finally, we connect
the construction to minimum metric realizations, metric spanners, and the
tight span from metric geometry.

Geodesic MDS is closely related to, but distinct from, several existing
methods.  \emph{Isomap} \citep{tenenbaum2000global} uses graph geodesics
on the input side only.  \emph{Kamada-Kawai graph drawing}
\citep{kamada1989algorithm} minimizes a stress between graph-theoretic
shortest-path distances and Euclidean drawing distances, but is restricted
to 2D visualization and is not formulated as a dimensionality reduction
method with explicit input/output distance comparison.  \emph{Metric graph
embedding} theory \citep{bourgain1985lipschitz,
linial1994geometry} studies distortion bounds for embedding graph metrics
into $\ell_p$ spaces but does not define a stress function or optimization
criterion.  \emph{Diffusion maps} \citep{coifman2006diffusion} define a
Markov-chain-based distance that is intrinsic to the data graph but again
uses Euclidean output distances.  To our knowledge, the self-consistent
formulation proposed here, with graph geodesics on the input side and
embedded lengths of the corresponding fixed input geodesics on the output
side, has not been previously proposed.

The remainder of the paper is organized as follows.
\cref{sec:background} reviews MDS, Isomap, and the $k$-NN graph construction.
\cref{sec:geodesic_mds} presents Geodesic MDS: the formal definition of
geodesic stress, the gradient derivation, the SMACOF extension, and the
geometric properties of the formulation.
Subsequent sections (not included in this draft) present experiments,
theoretical analysis, and applications.

%% ============================================================
\section{Background}
\label{sec:background}

Let $\Delta = (\delta_{ij})$ be an $n \times n$ symmetric dissimilarity
matrix with $\delta_{ii} = 0$.  Classical MDS
\citep{torgerson1952multidimensional, young1938discussion} seeks a
configuration $\bZ \in \R^{n \times d}$ such that
$\norm{\bz_i - \bz_j} \approx \delta_{ij}$.  When $\Delta$ contains
squared Euclidean distances, the solution is obtained by eigendecomposition
of the \emph{double-centered} Gram matrix
\begin{equation}
  \bB = -\tfrac{1}{2} \bJ \bD^{(2)} \bJ,
  \label{eq:gram}
\end{equation}
where $\bD^{(2)}_{ij} = \delta_{ij}^2$ and
$\bJ = \bI_n - n^{-1}\mathbf{1}\mathbf{1}^\top$ is the centering matrix.
The $d$-dimensional embedding is
$\bZ = \bU_d \Lambda_d^{1/2}$, where $\bU_d$ and $\Lambda_d$ contain the
$d$ largest eigenvectors and eigenvalues of $\bB$
\citep{eckart1936approximation, gower1966some}.
This is also known as \emph{principal coordinate analysis} (PCoA)
\citep{gower1966some}.  Classical MDS is a closed-form, non-iterative
method; it yields an exact isometric embedding when $\bB$ is positive
semidefinite (equivalently, when $\Delta$ satisfies Euclidean metric
conditions \citep{schoenberg1935remarks}).

When $\Delta$ is not a Euclidean distance matrix, classical MDS may place
negative eigenvalues in $\bB$, and the embedding is approximate.  Metric
MDS \citep{kruskal1964multidimensional} frames the problem explicitly as
an optimization: find $\bZ \in \R^{n \times d}$ minimizing a \emph{stress}
function that measures discrepancy between embedded distances and input
dissimilarities.

The most widely used criterion is Kruskal's \emph{Stress-1}
\citep{kruskal1964multidimensional}:
\begin{equation}
  \StressOne(\bZ) = \sqrt{
    \frac{\displaystyle\sum_{i < j}
          \bigl(d_{ij}(\bZ) - \delta_{ij}\bigr)^2}
         {\displaystyle\sum_{i < j} d_{ij}(\bZ)^2}
  },
  \label{eq:stress1}
\end{equation}
where $d_{ij}(\bZ) = \norm{\bz_i - \bz_j}$ is the Euclidean distance
between embedded points.  The normalization by $\sum_{i<j} d_{ij}^2$
makes $\StressOne$ invariant to uniform scaling of $\bZ$; it takes values
in $[0, 1]$, with 0 indicating a perfect isometric embedding.  An
alternative normalization divides by $\sum_{i<j} \delta_{ij}^2$
\citep{borg2005modern}, and is sometimes called Stress-1r; both are in
common use.

A generalization incorporates pair-specific weights $w_{ij} \geq 0$:
\begin{equation}
  \sigma_w(\bZ)
  = \sum_{i < j} w_{ij}\,\bigl(d_{ij}(\bZ) - \delta_{ij}\bigr)^2.
  \label{eq:weighted_stress}
\end{equation}
Common choices are $w_{ij} = \delta_{ij}^{-2}$ (relative weighting,
emphasizing small distances) and $w_{ij} = 1$ (uniform).  The unweighted
raw stress $\sigma(\bZ) = \sum_{i<j}(d_{ij}(\bZ) - \delta_{ij})^2$
is the numerator of $\StressOne^2$ up to scaling.

Empirical guidelines due to \citet{kruskal1964nonmetric} and
\citet{borg2005modern} suggest that $\StressOne < 0.05$ indicates an
excellent fit, $0.05$--$0.10$ a good fit, $0.10$--$0.20$ a fair fit, and
$> 0.20$ a poor fit, though these thresholds depend on $n$ and the
structure of $\Delta$.

Direct minimization of $\sigma$ or $\StressOne$ is a non-convex problem
with many local minima.  The \emph{Scaling by MAjorizing a COmplicated
Function} (SMACOF) algorithm \citep{deleeuw1977applications,
deleeuw2009multidimensional} is the dominant approach.  SMACOF exploits
the identity
\begin{equation}
  \sigma(\bZ) = \sum_{i<j}w_{ij}\delta_{ij}^2
               - 2\,\mathrm{tr}(\bZ^\top \bV\,\bZ)
               + \mathrm{tr}(\bZ^\top \bW\,\bZ),
  \label{eq:smacof_decomp}
\end{equation}
where $\bW_{ii} = \sum_{j \neq i} w_{ij}$, $\bW_{ij} = -w_{ij}$ ($i \neq
j$), and $\bV$ is the Guttman transform of $\bZ$:
\begin{equation}
  \bV(\bZ)_{ij} = -w_{ij}\,\frac{\delta_{ij}}{d_{ij}(\bZ)}\quad(i \neq j),
  \qquad
  \bV(\bZ)_{ii} = -\sum_{j \neq i} \bV(\bZ)_{ij}.
  \label{eq:guttman}
\end{equation}
The majorization step replaces the complicated function of $d_{ij}(\bZ)$
by a quadratic upper bound that is tight at the current iterate $\bZ^{(t)}$,
yielding the closed-form Guttman update
\begin{equation}
  \bZ^{(t+1)} = \bW^+ \bV(\bZ^{(t)})\,\bZ^{(t)},
  \label{eq:guttman_update}
\end{equation}
where $\bW^+$ is the Moore-Penrose pseudoinverse of $\bW$.  SMACOF
guarantees monotone decrease of $\sigma$ at each iteration and converges
to a stationary point \citep{deleeuw2009multidimensional}.  Multiple random
restarts are used in practice to mitigate local minima.  Efficient
implementations are available in R \citep{deleeuw2009multidimensional}
and Python \citep{pedregosa2011scikit}.

Non-metric MDS \citep{kruskal1964nonmetric, shepard1962analysis} relaxes
the requirement to reproduce input dissimilarity values, requiring only
that the \emph{rank order} of distances be preserved.  The optimization
alternates between isotonic regression (computing optimal
\emph{disparities} $\hat\delta_{ij}$ that are monotone in $\delta_{ij}$)
and gradient steps on the embedding.  Non-metric MDS is appropriate for
ordinal or perceptual data; Geodesic MDS as proposed here is a metric
method, though a non-metric geodesic variant is straightforward to
formulate by applying isotonic regression to the graph-geodesic targets
$d_G(i,j)$ while continuing to evaluate the output configuration by the
corresponding embedded path lengths $\demb_G(i,j;\bZ)$.

\emph{Isomap} \citep{tenenbaum2000global} addresses the limitation of MDS
on structured data whose relevant geometry is not well captured by ambient Euclidean distance by replacing Euclidean input
dissimilarities with approximations of intrinsic geodesic distances.
Given $\bX \in \R^{n \times m}$, Isomap first constructs a weighted
undirected $k$-NN graph $\Gk$ with vertex set $\{1, \ldots, n\}$,
connecting $i$ and $j$ by an edge of weight $\norm{\bx_i - \bx_j}_{\R^m}$
whenever $j$ is among the $k$ nearest neighbors of $i$ (or vice versa, in
the symmetric variant).  It then computes all-pairs shortest-path
distances $\dknn(i,j)$ on $\Gk$ using Dijkstra's algorithm
\citep{dijkstra1959note} or Floyd-Warshall \citep{floyd1962algorithm},
and sets $\delta_{ij} := \dknn(i,j)$.  Finally, it applies classical MDS
to the resulting dissimilarity matrix $\Delta$ to obtain the
low-dimensional embedding $\bZ \in \R^{n \times d}$.

In the special case of data sampled i.i.d.\ from a compact Riemannian manifold $(M, g)$
embedded in $\R^m$, the graph-geodesic distance $\dknn(i,j)$ converges
to the Riemannian geodesic distance $d_M(i,j)$ as $n \to \infty$, provided
$k$ grows appropriately with $n$ \citep{bernstein2000graph,
tenenbaum2000global, zha2003continuum, donoho2003hessian}.
The resulting classical MDS embedding then recovers the isometric embedding
of $M$ into a Euclidean space (when one exists) \citep{tenenbaum2000global}.

Despite this theoretical guarantee, Isomap measures the fidelity of the
embedding by Euclidean distances $\norm{\bz_i - \bz_j}$ between embedded
points, not by path-length distances along the embedded graph or geometric
scaffold.  For data whose intrinsic geometry is smooth and locally nearly
flat, this is a minor issue: the embedding is locally almost Euclidean, so
Euclidean and geodesic distances
are approximately equal.  For discrete graph metrics, branching metrics
(trees, stars), or data with topological bottlenecks, the discrepancy is
fundamental, as we elaborate in \cref{sec:geodesic_mds}.

The $k$-NN graph $\Gk$ arising in Isomap can be understood
as a \emph{metric graph}, a weighted graph whose path metric approximates
some underlying geodesic metric space.  Several observations are relevant
for our construction.

The edge weights of $\Gk$ are Euclidean distances in $\R^m$, so the path
metric $d_{\Gk}$ always satisfies $d_{\Gk}(i,j) \geq \norm{\bx_i - \bx_j}$
(the shortest path cannot be shorter than the straight line).  Equality
holds only for pairs connected by a direct edge.

$\Gk$ has at most $kn$ edges (exactly $kn/2$ in expectation for large $n$
with a symmetric construction), so the graph is sparse for small $k$.
As $k$ increases, more edges are added, $d_{\Gk}(i,j)$ decreases toward
$\norm{\bx_i - \bx_j}$, and the graph metric approaches the Euclidean
ambient metric.  At $k = n-1$ (complete graph), $d_{\Gk}(i,j) =
\norm{\bx_i - \bx_j}$ for all pairs, and the construction degenerates to
standard Euclidean MDS.

$\Gk$ is a $(1+\varepsilon)$-spanner of the Euclidean metric on $\bX$ for
sufficiently large $k$ depending on the intrinsic dimensionality of the
data \citep{althofer1993sparse, har-peled2006low}.  This means
$d_{\Gk}(i,j) \leq (1+\varepsilon)\norm{\bx_i - \bx_j}$, so the graph
geodesic approximates the Euclidean distance with controlled relative
error.  For data with low intrinsic dimension $d^*$, $k = O(\varepsilon^{-d^*})$
suffices \citep{har-peled2006low}.

An edge $(i,j)$ in $K_n$ (the complete graph on $\bX$) is
\emph{metrically redundant} with respect to $d_{\Gk}$ if there exists a
path through intermediate points with equal total length.  The subgraph
obtained by removing all metrically redundant edges is the
\emph{minimum realization} $G^*$ of the metric $d_{\Gk}$
\citep{dress1984trees, imrich1984optimal}, and encodes the intrinsic
geodesic structure of the data.  For tree metrics, $G^*$ is the unique
spanning tree; for hyperbolic metrics it is sparse; for generic metrics
it approaches the complete graph.
Figure~\ref{fig:k_regimes} is a natural companion here because the role of
$k$ is one of the main geometric ideas of the paper.

\begin{figure}[t]
  \centering
  \includegraphics[width=0.98\linewidth]{../figures/fig_k_regimes.pdf}
  \caption{The neighborhood parameter $k$ acts as a geometric scale parameter.  The same horseshoe-shaped point set is shown for three neighborhood sizes, with the same endpoint pair $p,q$ highlighted in each panel.  For small $k$, the graph is sparse and nearly tree-like; for intermediate $k$, shortest paths still follow the underlying geometry and preserve local geodesic structure; and for large $k$, dense shortcut edges appear and the metric approaches the ambient Euclidean regime.}
  \label{fig:k_regimes}
\end{figure}

The broader context of our work is the literature often grouped under
\emph{manifold learning}.  Many real datasets are not well modeled as
samples from a manifold in the strict sense, but the common goal is still
to find low-dimensional representations of high-dimensional data with
non-Euclidean local or shortest-path structure.  Beyond Isomap, influential methods
include Locally Linear Embedding (LLE) \citep{roweis2000nonlinear},
Laplacian Eigenmaps \citep{belkin2003laplacian}, Hessian Eigenmaps
\citep{donoho2003hessian}, Diffusion Maps \citep{coifman2006diffusion},
$t$-SNE \citep{vandermaaten2008tsne}, and UMAP \citep{mcinnes2018umap}.
With the exception of diffusion maps (which define an intrinsic distance via
Markov chain commute times), all of these methods use Euclidean distances
on the output side, making the inconsistency identified here relevant to the
broader nonlinear-embedding literature.

Graph drawing methods, particularly Kamada-Kawai \citep{kamada1989algorithm}
and Fruchterman-Reingold \citep{fruchterman1991graph}, are also related.
Kamada-Kawai minimizes a stress between graph-theoretic shortest-path
distances (in the unweighted graph) and Euclidean distances in the 2D
drawing, which is structurally similar to Isomap applied to an unweighted
graph.  A geodesic variant of Kamada-Kawai, using embedded lengths of
fixed graph geodesics on the output, is a natural companion to Geodesic
MDS but is outside the scope of the present paper.

%% ============================================================
\section{Geodesic MDS}
\label{sec:geodesic_mds}

\subsection{The Self-Consistency Principle}
\label{sec:self_consistency}

We begin by articulating the principle that motivates the proposed method.

We call an embedding $f: X \to \R^d$ \emph{self-consistent} with respect
to $(G, \Gamma)$ when $(X,\delta)$ is a metric space with graph
structure $G = (X, E, w)$ whose path metric $d_G$ satisfies
$d_G(i,j) \geq \delta(i,j)$ for all $i,j$, a canonical $d_G$-geodesic
$\gamma_{ij}$ is chosen for each pair $(i,j)$, and the fidelity of the
embedding is measured by the path length of the image $f(\gamma_{ij})$,
that is, by the sum of Euclidean edge lengths along that same input-side
geodesic.

Standard metric MDS and Isomap are not self-consistent in this sense:
they measure fidelity by Euclidean distances $\norm{f(i) - f(j)}$, which
are at most equal to any path-length distance through intermediate vertices.

The self-consistency principle has a natural interpretation:
\emph{the distance used to evaluate the embedding should be the same
type of distance, along the same family of geodesics, used to construct
the input dissimilarities.}
If input dissimilarities are graph geodesics, output distances should be
measured along those same graph geodesics after embedding.  This is not
merely a philosophical preference: the
analysis in \cref{sec:systematic_bias} below shows that using Euclidean
output distances introduces a systematic downward bias in the estimated
distances for topologically constrained pairs.

\subsection{Systematic Bias of Euclidean Output Distances}
\label{sec:systematic_bias}

The output-side bias follows from a simple triangle-inequality
observation.  Let $\bZ = (\bz_1, \ldots, \bz_n) \in \R^{n \times d}$ be
any embedding, and let $G = (V, E, w)$ be a weighted graph with vertex
set $V = \{1, \ldots, n\}$.  For any pair $(i,j)$ with chosen shortest
path $i = v_0, v_1, \ldots, v_\ell = j$ in $G$, define the embedded
length of that geodesic by
\begin{equation}
  \demb_G(i,j;\bZ)
  = \sum_{s=0}^{\ell-1} \norm{\bz_{v_s} - \bz_{v_{s+1}}}.
  \label{eq:emb_geo}
\end{equation}
Then
\begin{equation}
  \norm{\bz_i - \bz_j} \;\leq\; \demb_G(i,j;\bZ),
  \label{eq:triangle_bias}
\end{equation}
with equality if and only if all embedded path vertices
$\bz_{v_0}, \ldots, \bz_{v_\ell}$ are collinear in $\R^d$ in the order
of the path.  Indeed, by the triangle inequality applied $\ell$ times,
$\norm{\bz_i - \bz_j} = \norm{\sum_{s} (\bz_{v_{s+1}} - \bz_{v_s})}
\leq \sum_s \norm{\bz_{v_{s+1}} - \bz_{v_s}} = \demb_G(i,j;\bZ)$, and
equality holds precisely when all difference vectors
$\bz_{v_{s+1}} - \bz_{v_s}$ are non-negative scalar multiples of a single
unit vector.

This triangle-inequality observation shows that standard MDS and Isomap systematically
underestimate embedded geodesic lengths for all topologically
non-degenerate pairs (those whose shortest path has $\ell \geq 2$ and
does not embed as a straight line).  When the input dissimilarity
$\delta_{ij}$ is a graph geodesic, minimizing Euclidean-output stress
creates pressure to satisfy $\norm{\bz_i - \bz_j} \approx \delta_{ij}$,
which in turn forces embedded paths to be approximately collinear.
For tree-structured data this collapses the tree into a lower-dimensional
structure than warranted; for star graphs it forces nearly collinear arm
arrangement regardless of the dimension $d$.

\subsection{The Geodesic Stress Function}
\label{sec:geodesic_stress}

We now define the proposed method.

\subsubsection*{Setup}
Let $G = (V,E,w)$ be a connected weighted graph with vertex set
$V = \{1,\ldots,n\}$ and positive edge weights $w_{uv} > 0$.  Write
$d_G$ for the weighted shortest-path metric on $V$.  For each pair
$i < j$, choose a canonical $d_G$-geodesic
\begin{equation}
  \gamma_{ij} = (v_0 = i, v_1, \ldots, v_L = j),
  \label{eq:canonical_geodesic}
\end{equation}
breaking ties deterministically when multiple shortest paths exist.  The
collection $\Gamma = \{\gamma_{ij} : i < j\}$ is part of the data of the
method.  In this sense, Geodesic MDS is defined on the pair $(G,\Gamma)$,
not on $G$ alone.

An embedding is a configuration matrix $\bZ = (\bz_1, \ldots, \bz_n)^\top
\in \R^{n \times d}$.  For fixed $(G,\Gamma)$, the \emph{embedded graph}
$G(\bZ)$ is the same graph $G$ with edge lengths replaced by embedded
Euclidean distances.  The \emph{embedded length of the chosen input
geodesic} is
\begin{equation}
  \demb_{\Gamma}(i,j;\bZ)
  = \sum_{s=0}^{L-1} \norm{\bz_{v_s} - \bz_{v_{s+1}}},
  \label{eq:emb_geodesic}
\end{equation}
where $\gamma_{ij} = (v_0, \ldots, v_L)$.

We define the \emph{raw geodesic stress} of an embedding
$\bZ \in \R^{n \times d}$ with respect to $(G,\Gamma)$ by
\begin{equation}
  \Stress^{\mathrm{geo}}_{\Gamma}(\bZ; G)
  = \sum_{i < j}
    \bigl(\demb_{\Gamma}(i,j;\bZ) - d_G(i,j)\bigr)^2.
  \label{eq:raw_geo_stress}
\end{equation}

The normalized \emph{geodesic Stress-1} of $\bZ$ with respect to
$(G,\Gamma)$ is
\begin{equation}
  {\StressOne}_{\Gamma}^{\mathrm{geo}}(\bZ; G)
  = \sqrt{
      \frac{\Stress^{\mathrm{geo}}_{\Gamma}(\bZ; G)}
           {\displaystyle\sum_{i < j} d_G(i,j)^2}
    }.
  \label{eq:geo_stress}
\end{equation}

A \emph{Geodesic MDS embedding} of $(G,\Gamma)$ in $\R^d$ is any
minimizer
\begin{equation}
  \hat{\bZ}_{\Gamma,d}(G)
  \in \argmin_{\bZ \in \R^{n \times d}}
  \Stress^{\mathrm{geo}}_{\Gamma}(\bZ; G).
  \label{eq:geo_mds}
\end{equation}
Equivalently, since the denominator in \cref{eq:geo_stress} is constant
in $\bZ$, one may minimize ${\StressOne}_{\Gamma}^{\mathrm{geo}}(\bZ; G)$.

\subsubsection*{Specialization to $k$-NN graphs}
When the input is a data matrix $\bX \in \R^{n \times m}$, we recover the
construction used throughout this paper by fixing a neighborhood parameter
$k \in \{1,\ldots,n-1\}$ and building the symmetric $k$-NN graph
$\Gk = (V, E_k)$ with
\begin{equation}
  E_k = \{(i,j) : j \in \mathcal{N}_k(i) \text{ or } i \in \mathcal{N}_k(j)\},
  \label{eq:knn_edges}
\end{equation}
where $\mathcal{N}_k(i)$ denotes the $k$ nearest neighbors of $i$ in
$\R^m$.  Assign weights $w_{ij} = \norm{\bx_i - \bx_j}_{\R^m}$ for
$(i,j) \in E_k$, compute the shortest-path metric
$d_{\Gk}(i,j) = \delta_{ij}^{(k)}$, and choose a canonical shortest-path
family $\Gamma_k = \{\gamma_{ij}^{(k)} : i < j\}$.  The resulting Geodesic
MDS objective is
\begin{equation}
  \Stress^{\mathrm{geo}}_{\Gamma_k}(\bZ; \Gk)
  = \sum_{i < j}
    \bigl(\demb_{\Gamma_k}(i,j;\bZ) - \delta^{(k)}_{ij}\bigr)^2.
  \label{eq:knn_raw_geo_stress}
\end{equation}

We normalize by the constant $\sum_{i<j} d_G(i,j)^2$ rather than by an
embedded denominator such as $\sum_{i<j} \demb_{\Gamma}(i,j;\bZ)^2$.
This keeps ${\StressOne}_{\Gamma}^{\mathrm{geo}}$ scale-free while
leaving the optimizer identical to the raw numerator
\cref{eq:raw_geo_stress}.  A Kruskal-style denominator depending on
$\bZ$ is still useful for reporting, but it changes the optimization
target.

The graph topology $E$ and the chosen geodesic family $\Gamma$ are fixed
by the input metric data and are not part of the optimization.  Only the
vertex positions $\bz_1,\ldots,\bz_n$ are free.  This is a key design
choice: it ensures that the relevant routes are determined by the input
geometry and that the optimization preserves that chosen geodesic
structure after embedding.
Figure~\ref{fig:geodesic_stress_definition} should make the fixed-geodesic,
two-metric comparison in \cref{eq:emb_geodesic,eq:geo_stress} easy to see.

\begin{figure}[t]
  \centering
  \includegraphics[width=0.98\linewidth]{../figures/fig_geodesic_stress_definition.pdf}
  \caption{Definition of geodesic stress.  Left: the input
  dissimilarity $d_G(i,j)$ is the shortest-path distance between $i$ and
  $j$ on the fixed weighted graph $G$; in the data-analytic specialization,
  $G = \Gk$ is a $k$-NN graph.  Right: the output quantity
  $\demb_{\Gamma}(i,j;\bZ)$ is the length of the canonical shortest path
  chosen in the input graph, measured after the vertex positions move in
  $\R^d$.  The route is fixed by the input geometry; only the embedded edge
  lengths change during optimization.}
  \label{fig:geodesic_stress_definition}
\end{figure}


The objective \cref{eq:raw_geo_stress} is a fidelity criterion, not an
aesthetics criterion, so its minimizers need not be unique.  This is a
feature of the definition: multiple embeddings can realize the same graph
geodesic data equally well.  For example, for a path graph
$v_0 - v_1 - v_2$ with any fixed edge lengths,
every embedding in $\R^2$ that preserves those two edge lengths is a global
minimum of \cref{eq:raw_geo_stress}; one obtains an infinite family of minima
by fixing the segment $[v_0,v_1]$ and rotating $[v_1,v_2]$ around $v_1$.
Likewise, if $v_0,v_1,v_2$ form one triangle and $v_1,v_2,v_3$ form another
triangle sharing the edge $[v_1,v_2]$, then both the non-overlapping and the
overlapping realizations are global minima whenever all edge lengths are
preserved.  \Cref{fig:nonuniqueness_examples} illustrates these two basic
sources of nonuniqueness.

\begin{figure}[t]
  \centering
  \includegraphics[width=0.98\linewidth]{../figures/fig_nonuniqueness_examples.pdf}
  \caption{Two sources of nonuniqueness for geodesic-stress minimizers.
  Left: a three-vertex path has a one-parameter family of exact realizations
  obtained by rotating the second edge while preserving both edge lengths.
  Right: two triangles sharing an edge admit distinct exact realizations,
  including a non-overlapping embedding and an overlapping embedding.  In both
  examples the raw geodesic stress is zero because all relevant graph-geodesic
  lengths are preserved.}
  \label{fig:nonuniqueness_examples}
\end{figure}

\subsubsection*{A diagnostic lesson}
The diagnostic experiments reported separately in this project show that the
all-pairs objective \cref{eq:raw_geo_stress} can admit visually pathological
minimizers on simple synthetic graphs.  On orthogonal grids, graph geodesics
systematically exceed straight-line Euclidean distances, so the gradient of the
all-pairs geodesic stress can push the embedding toward stretched or folded
configurations.  This does not invalidate the definition of GMDS: it shows that
regularization is a separate modeling choice that may be needed when one wants
both fidelity to graph geodesics and a visually well-spread realization.

\subsubsection*{Regularized GMDS}
To incorporate such modeling choices explicitly, introduce a Euclidean
dispersion regularizer
\begin{equation}
  R(\bZ) = \sum_{i < j} f_r\bigl(\norm{\bz_i - \bz_j}\bigr),
  \label{eq:repulsive_regularizer}
\end{equation}
where $f_r : (0,\infty) \to (0,\infty)$ is a decreasing function.  A
concrete choice used widely in force-directed graph drawing is the Coulomb-type
penalty
\begin{equation}
  f_r(r) = \frac{1}{r^2 + \varepsilon_r^2},
  \label{eq:coulomb_regularizer}
\end{equation}
with a small $\varepsilon_r > 0$ to prevent singularities.  The
self-consistency principle applies to the fidelity term
\cref{eq:raw_geo_stress}; the role of $R$ is different.  It is a Euclidean
regularizer that encourages spread and discourages collapse, not a substitute
for the graph-geodesic mismatch term itself.  A standard scalarized objective
is therefore
\begin{equation}
  F_\lambda(\bZ; G,\Gamma)
  = \Stress^{\mathrm{geo}}_{\Gamma}(\bZ; G) + \lambda R(\bZ),
  \qquad \lambda > 0.
  \label{eq:regularized_gmds}
\end{equation}
A \emph{regularized GMDS embedding} is a local minimizer of
\cref{eq:regularized_gmds} for a specified $\lambda$.  The value of
$\lambda$ is not universal; in practice it should be reported explicitly and
chosen by a fidelity-dispersion trade-off study, continuation scheme, or
task-specific validation criterion.

\subsubsection*{Geodesic KK, LGKK, and edge-based relaxations}
A useful bridge to graph drawing is obtained by specializing the GMDS
criterion to Kamada-Kawai-style pair weights.  With a global stiffness
constant $K_0 > 0$, pairwise graph distances $d_G(i,j)$, and a fitted
scale $L_0 > 0$, define the geodesic Kamada-Kawai objective
\begin{equation}
  E_{\Gamma,L_0}^{\mathrm{gKK}}(\bZ)
  = \tfrac{1}{2}\sum_{i<j}
    k_{ij}\bigl(\demb_{\Gamma}(i,j;\bZ) - L_0\,d_G(i,j)\bigr)^2,
  \qquad
  k_{ij} = \frac{K_0}{\max\{d_G(i,j), \varepsilon_d\}^2},
  \label{eq:geodesic_kk}
\end{equation}
with a small $\varepsilon_d > 0$.  This is still a fixed-geodesic
geodesic-isometry criterion: it compares the input graph distances to the
embedded lengths of those same chosen input-side geodesics, but with
Kamada-Kawai-style stiffness weights and a fitted global scale.  In
\texttt{grip}, the full all-pairs version is implemented as geodesic KK,
and a sparse landmark restriction of the same objective is implemented as
landmark-geodesic KK (LGKK).

Classical Kamada-Kawai is recovered by replacing the embedded geodesic
lengths $\demb_{\Gamma}(i,j;\bZ)$ by straight-line Euclidean distances
$\norm{\bz_i - \bz_j}$.  In this sense, geodesic KK and LGKK form an
intermediate bridge between graph-drawing objectives and the more
principled GMDS criterion: they retain the input-side geodesic structure
of GMDS while simplifying the weighting and, in the LGKK case, the active
pair set.

If one simplifies further by restricting the sum to graph edges only, one
obtains the edge-only functional
\begin{equation}
  E_K(\bZ)
  = \sum_{(i,j) \in E} k_{ij}\bigl(\norm{\bz_i - \bz_j} - w_{ij}\bigr)^2,
  \label{eq:edge_relaxation}
\end{equation}
because for an edge $(i,j)$ the chosen graph geodesic is the edge itself and
thus $\demb_{\Gamma}(i,j;\bZ) = \norm{\bz_i - \bz_j}$.  This is no longer the
full geodesic KK or GMDS criterion: it penalizes only single-hop
distortion, and is therefore closer to an edge-restricted classical
Kamada-Kawai objective \citep{kamada1989algorithm, fruchterman1991graph}.

The relation among \cref{eq:raw_geo_stress,eq:geodesic_kk,eq:edge_relaxation}
is easiest to see in edge-length coordinates.  Let $\ell(\bZ) \in \R^{|E|}$
be the vector of embedded edge lengths, let $w \in \R^{|E|}$ collect the
graph edge weights, let $d \in \R^{\binom{n}{2}}$ collect the all-pairs
graph distances, and let $A_{\Gamma}$ be the path-edge incidence matrix
whose $(i,j)$-row indicates which edges lie on the chosen geodesic
$\gamma_{ij}$.  Then raw GMDS compares $A_{\Gamma}\ell(\bZ)$ to $d$ with
uniform all-pairs weights, while geodesic KK takes the weighted form
\begin{equation}
  E_{\Gamma,L_0}^{\mathrm{gKK}}(\bZ)
  = \tfrac{1}{2}\bigl(A_{\Gamma}\ell(\bZ) - L_0 d\bigr)^\top
    K_{\mathrm{pair}}
    \bigl(A_{\Gamma}\ell(\bZ) - L_0 d\bigr),
  \label{eq:path_incidence_gkk}
\end{equation}
with $K_{\mathrm{pair}} = \diag(k_{ij})$ diagonal over unordered pairs,
whereas the edge-only relaxation takes the form
\begin{equation}
  E_K(\bZ)
  = \bigl(\ell(\bZ)-w\bigr)^\top K_{\mathrm{edge}} \bigl(\ell(\bZ)-w\bigr),
  \label{eq:edge_relaxation_matrix}
\end{equation}
with $K_{\mathrm{edge}} = \diag(k_{ij})$ diagonal over edges.  Thus the
passage from GMDS to geodesic KK changes the pair weighting and introduces
a fitted global scale, whereas the passage from geodesic KK to
\cref{eq:edge_relaxation} drops the multi-hop pair set and the associated
path-coupling structure encoded by $A_{\Gamma}$.  What is gained is a much
cheaper objective; what is lost is explicit control of long-range
path-length errors.

If $E_K(\bZ) = 0$, then every edge length is realized exactly, hence every
chosen graph-geodesic length is realized exactly as well, and
\cref{eq:raw_geo_stress} also vanishes.  In the approximate regime, however,
the two objectives are different and generally have different local minima.
For this reason we treat \cref{eq:edge_relaxation} as a relaxation of both
geodesic KK and GMDS, not as the definition of GMDS itself.  This
perspective suggests a natural continuation strategy:
classical KK $\to$ geodesic KK or LGKK $\to$ GMDS.  More generally, one can
imagine a filtration of active pair sets, beginning with local or landmark
pairs and expanding toward the full all-pairs set, as a homotopy from
cheap graph-drawing surrogates to the full geodesic-isometry objective.
We do not analyze such schedules here, but they provide a useful bridge
between graph drawing and GMDS.

In the full geodesic KK objective \cref{eq:geodesic_kk}, the natural
default is the classical Kamada-Kawai stiffness
$k_{ij} = K_0/\max\{d_G(i,j), \varepsilon_d\}^2$.  In the edge-only
functional \cref{eq:edge_relaxation}, the weights $k_{ij}$ are
user-chosen tuning parameters.  A simple default is $k_{ij} \equiv 1$.
If one wants to emphasize edge lengths that occur frequently in the graph,
one may set $k_{ij} = \hat f_{\mathrm{edge}}(w_{ij})$, where
$\hat f_{\mathrm{edge}}$ is a kernel density estimate of the empirical
distribution of edge lengths.  Other choices, such as weights proportional to
geodesic-usage counts, are also possible and should be justified by the
application.

\subsection{Properties of Geodesic MDS}
\label{sec:properties}

Geodesic MDS is self-consistent by construction:
both input dissimilarities and output distances are measured on the same
fixed family of geodesics in $(G,\Gamma)$ (the former by the graph
metric $d_G$, the latter by the embedded path lengths
$\demb_{\Gamma}(i,j;\bZ)$).

When the input graph is the $k$-NN specialization and $k = n-1$, every
pair is connected by a direct edge, so
$\demb_{\Gamma_{n-1}}(i,j;\bZ) = \norm{\bz_i - \bz_j}$ and
$\delta^{(n-1)}_{ij} = \norm{\bx_i - \bx_j}$.  Geodesic stress reduces to
standard metric Stress-1, and Geodesic MDS recovers standard metric MDS.

For $k = 1$, $\Gk$ is (generically) a sparse forest; with a suitable
connectivity correction (e.g., taking the minimum spanning tree as
$G_1$) the graph is a tree.  The tree-path metric is then a tree metric in
the sense of \citet{dress1984trees}, and Geodesic MDS seeks an
embedding that faithfully represents the tree structure.  Tree metrics
embed isometrically into $\R^1$ (the real line) with zero stress;
thus for tree-structured data Geodesic MDS with $k=1$ achieves
perfect stress in $d=1$.

As $k$ increases from 1 to $n-1$, more edges are added to $\Gk$,
$\delta^{(k)}_{ij}$ decreases toward $\norm{\bx_i - \bx_j}$, and
$\demb_{\Gamma_k}(i,j;\bZ)$ decreases toward $\norm{\bz_i - \bz_j}$.  The
geodesic stress thus interpolates between a tree-metric embedding at small
$k$ and standard MDS at large $k$.  The parameter $k$ controls the
geometric scale at which geodesic structure is enforced, analogous to the
bandwidth parameter in kernel methods.

By the triangle-inequality observation above,
$\demb_{\Gamma_k}(i,j;\bZ) \geq \norm{\bz_i - \bz_j}$
for all $i,j$ and all $\bZ$.  If in addition $\delta^{(k)}_{ij} \geq
\norm{\bx_i - \bx_j}$ (which holds by the triangle inequality on graph
paths), then for the same configuration $\bZ$ the geodesic stress is
larger than the Euclidean stress.  The optimal geodesic configuration
$\hat{\bZ}_{\Gamma_k,d}(\Gk)$ is not in general the same as the optimal
Euclidean configuration.

\subsection{Gradient Derivation}
\label{sec:gradient}

We derive the gradient of the raw geodesic stress
$\Stress^{\mathrm{geo}}_{\Gamma_k}(\bZ; \Gk)$
(\cref{eq:knn_raw_geo_stress}) with respect to the position
$\bz_\ell \in \R^d$ of a single vertex $\ell$.  For brevity, write
$\sigma^{\mathrm{geo}}(\bZ) = \Stress^{\mathrm{geo}}_{\Gamma_k}(\bZ; \Gk)$.

\subsubsection*{Fixed geodesics and differentiability}
For each pair $i < j$, let
$\gamma_{ij}^{(k)} = (v_0 = i, v_1, \ldots, v_L = j)$ be the canonical
shortest path chosen once from the input graph $\Gk$.  Then
$\demb_{\Gamma_k}(i,j;\bZ)$ is the length of this fixed polygonal chain in the
embedding.  Since the route does not change with $\bZ$, the only
nonsmoothness comes from Euclidean norms at zero-length edges; away from
such degenerate configurations, $\demb_{\Gamma_k}$ is differentiable.

\subsubsection*{Gradient with respect to an intermediate vertex}
Let $\gamma_{ij}^{(k)} = (v_0, v_1, \ldots, v_L)$ be the chosen shortest
path from $i = v_0$ to $j = v_L$.  Define the unit edge vector
\begin{equation}
  \mathbf{e}_{ab} = \frac{\bz_a - \bz_b}{\norm{\bz_a - \bz_b}},
  \label{eq:unit_edge}
\end{equation}
so that $\frac{\partial}{\partial \bz_a}\norm{\bz_a - \bz_b} = \mathbf{e}_{ab}$
and $\frac{\partial}{\partial \bz_b}\norm{\bz_a - \bz_b} = -\mathbf{e}_{ab}$.

The gradient of $\demb_{\Gamma_k}(i,j;\bZ)$ with respect to $\bz_\ell$ depends on
whether $\ell$ appears in $\gamma_{ij}^{(k)}$:
\begin{align}
  \frac{\partial\, \demb_{\Gamma_k}(i,j;\bZ)}{\partial\, \bz_\ell}
  &= \begin{cases}
       \mathbf{e}_{\ell, v_{s-1}} + \mathbf{e}_{\ell, v_{s+1}}
         & \text{if } \ell = v_s,\ 0 < s < L, \\
       \mathbf{e}_{\ell, v_1}
         & \text{if } \ell = v_0 = i, \\
       \mathbf{e}_{\ell, v_{L-1}}
         & \text{if } \ell = v_L = j, \\
       \bm{0}
         & \text{if } \ell \notin \gamma_{ij}^{(k)}.
     \end{cases}
  \label{eq:geo_distance_gradient}
\end{align}

The formula captures the non-local nature of geodesic MDS: only those
pairs $(i,j)$ whose chosen input geodesic passes through vertex $\ell$ contribute
to the gradient at $\bz_\ell$.
Figure~\ref{fig:gradient_nonlocality} should help readers see why the
gradient depends on global geodesic-route structure rather than only on
immediate graph neighbors.

\begin{figure}[t]
  \centering
  \begin{minipage}[t]{0.82\textwidth}
    \centering
    \includegraphics[width=\linewidth]{../figures/fig_gradient_nonlocality_paths.pdf}
  \end{minipage}

  \caption{Geodesic stress has non-local gradients.  A vertex receives contributions only from endpoint pairs whose chosen input geodesic passes through it, so the gradient depends on global route structure rather than only on immediate graph neighbors.}
  \label{fig:gradient_nonlocality}
\end{figure}

\subsubsection*{Gradient of the raw geodesic stress}
By the chain rule,
\begin{equation}
  \frac{\partial\, \sigma^{\mathrm{geo}}}{\partial\, \bz_\ell}
  = 2 \sum_{i < j}
    \bigl(\demb_{\Gamma_k}(i,j;\bZ) - \delta^{(k)}_{ij}\bigr)
    \cdot \frac{\partial\, \demb_{\Gamma_k}(i,j;\bZ)}{\partial\, \bz_\ell}.
  \label{eq:raw_gradient}
\end{equation}
Substituting \cref{eq:geo_distance_gradient} and collecting terms,
the gradient decomposes into contributions from paths that pass through
$\ell$ as an \emph{internal} vertex vs.\ paths that have $\ell$ as an
\emph{endpoint}:
\begin{align}
  \frac{\partial\, \sigma^{\mathrm{geo}}}{\partial\, \bz_\ell}
  &= 2 \sum_{\substack{i < j \\ \ell \in \gamma_{ij}^{(k)},\, \ell \neq i,j}}
     \!\!\!\!\!\!
     r_{ij}\,\bigl(\mathbf{e}_{\ell,\mathrm{prev}(\gamma_{ij}^{(k)})}
                   + \mathbf{e}_{\ell,\mathrm{next}(\gamma_{ij}^{(k)})}\bigr)
     \nonumber \\
  &\quad + 2\sum_{\substack{j > \ell \\ \ell \in \gamma_{\ell j}^{(k)}}}
     r_{\ell j}\,\mathbf{e}_{\ell,\mathrm{next}(\gamma_{\ell j}^{(k)})}
     + 2\sum_{\substack{i < \ell \\ \ell \in \gamma_{i\ell}^{(k)}}}
     r_{i\ell}\,\mathbf{e}_{\ell,\mathrm{prev}(\gamma_{i\ell}^{(k)})},
  \label{eq:full_gradient}
\end{align}
where $r_{ij} = \demb_{\Gamma_k}(i,j;\bZ) - \delta^{(k)}_{ij}$ is the signed
residual and $\mathrm{prev}(\gamma_{ij}^{(k)})$ (resp.\
$\mathrm{next}(\gamma_{ij}^{(k)})$) denotes the predecessor (resp.\
successor) of $\ell$ on the fixed path $\gamma_{ij}^{(k)}$.  The entire
gradient matrix $\nabla_{\bZ} \sigma^{\mathrm{geo}}
\in \R^{n \times d}$ can be assembled in $O(n^2 L_{\max})$ operations,
where $L_{\max}$ is the maximum path length in $\Gk$.

\subsubsection*{Degenerate edge lengths}
The geodesic stress is not differentiable only when some embedded edge on
a chosen geodesic collapses to zero length.  At such a point, one may use
the standard subgradient of the Euclidean norm or introduce a small
edge-length stabilizer.  Away from these measure-zero degeneracies, the
gradient above is valid.

\subsection{SMACOF Majorization for Geodesic Stress}
\label{sec:geo_smacof}

The standard SMACOF majorization exploits the representation
$d_{ij}(\bZ) = \norm{\bz_i - \bz_j}$ as a single Euclidean norm, which
can be bounded below by a linear function of $\bZ$ via the Cauchy-Schwarz
inequality.  For geodesic MDS, $\demb_{\Gamma_k}(i,j;\bZ)$ is a \emph{sum} of
Euclidean norms along a path.  This structure still admits a majorization,
as we now show.

\subsubsection*{Fixed-geodesic majorization}
Fix a reference configuration $\bZ^{(t)}$.  The path family $\Gamma_k$ is
already fixed by the input metric, so $\demb_{\Gamma_k}(i,j;\bZ)$ is always a sum
of norms along the stored path $\gamma_{ij}^{(k)}$.  Each
individual norm $\norm{\bz_a - \bz_b}$ admits the Cauchy-Schwarz majorization
\begin{equation}
  \norm{\bz_a - \bz_b}
  \leq
  \frac{\norm{\bz_a - \bz_b}^2 + \norm{\bz^{(t)}_a - \bz^{(t)}_b}^2}
       {2\,\norm{\bz^{(t)}_a - \bz^{(t)}_b}}
  \quad \text{(tight at } \bZ = \bZ^{(t)}\text{)}.
  \label{eq:cs_majorization}
\end{equation}
Summing over path edges, the total geodesic distance is majorized by a
function that is quadratic in $\bZ$, and hence the stress is majorized by
a fourth-degree polynomial.  A tighter approach uses the standard SMACOF
identity applied path-edge by path-edge.

\subsubsection*{Path-augmented Guttman transform}
Define the \emph{path-edge weight} $\tilde{w}_{ab}^{(t)}$ as the total
weight accumulated by edge $(a,b)$ across all chosen input geodesics:
\begin{equation}
  \tilde{w}_{ab}^{(t)}
  = \sum_{i < j : (a,b) \in \gamma_{ij}^{(k)}}
    \frac{1}{\norm{\bz^{(t)}_a - \bz^{(t)}_b}},
  \label{eq:path_edge_weight}
\end{equation}
and the \emph{path-augmented Guttman matrix}
\begin{equation}
  \tilde{\bV}^{(t)}_{ab}
  = -\tilde{w}^{(t)}_{ab}, \quad a \neq b; \qquad
  \tilde{\bV}^{(t)}_{aa}
  = \sum_{b \neq a} \tilde{w}^{(t)}_{ab}.
  \label{eq:augmented_guttman}
\end{equation}
The majorization step then takes the form
\begin{equation}
  \bZ^{(t+1)} = \bigl(\tilde{\bW}^{(t)}\bigr)^+ \tilde{\bV}^{(t)} \bZ^{(t)},
  \label{eq:geo_guttman_update}
\end{equation}
where $\tilde{\bW}^{(t)}$ is the diagonal matrix of row sums of
$|\tilde{\bV}^{(t)}|$.  Since the geodesic family is fixed, the stored
paths are reused at every iteration and only the embedded edge lengths are
updated.

With the geodesic family $\Gamma_k$ fixed by the input metric, each step
\eqref{eq:geo_guttman_update} does not increase
$\sigma^{\mathrm{geo}}(\bZ)$.  The reason is that at iteration $t$, the
fixed-geodesic majorization gives an upper bound
$\sigma^{\mathrm{geo}}(\bZ) \leq f^{(t)}(\bZ)$ for a function
$f^{(t)}$ that is quadratic in $\bZ$ and tight at $\bZ^{(t)}$.  The
update \eqref{eq:geo_guttman_update} minimizes $f^{(t)}$, so
$f^{(t)}(\bZ^{(t+1)}) \leq f^{(t)}(\bZ^{(t)}) =
\sigma^{\mathrm{geo}}(\bZ^{(t)})$.  Because the same fixed path family is
used in both the objective and the majorizer, we also have
$\sigma^{\mathrm{geo}}(\bZ^{(t+1)}) \leq f^{(t)}(\bZ^{(t+1)})$, hence
\[
  \sigma^{\mathrm{geo}}(\bZ^{(t+1)})
  \leq f^{(t)}(\bZ^{(t+1)})
  \leq f^{(t)}(\bZ^{(t)})
  = \sigma^{\mathrm{geo}}(\bZ^{(t)}).
\]

\subsubsection*{Computational complexity}
Each iteration of the geodesic SMACOF requires:
(i) path accumulation for $\tilde{\bW}^{(t)}$ and $\tilde{\bV}^{(t)}$
using the stored geodesic family $\Gamma_k$: $O(n^2 L_{\max})$;
(ii) a linear system solve: $O(n^3)$ or $O(n^2 d)$ with conjugate gradient.
The dominant cost is $O(n^3)$, the same as standard SMACOF.
The all-pairs shortest paths and the canonical geodesic family are
computed once in preprocessing at cost
$O(n(|E_k| + n\log n))$ using Dijkstra with a Fibonacci heap
\citep{fredman1987fibonacci}.
Figure~\ref{fig:geodesic_mds_workflow} summarizes the method at the level of
data flow and is intended to complement the pseudocode in
\cref{alg:geodesic_mds}.

\begin{figure}[t]
  \centering
  \includegraphics[width=0.98\linewidth]{../figures/fig_geodesic_mds_workflow.pdf}
  \caption{One iteration of Geodesic MDS.  A canonical family of shortest paths is chosen once from the input graph.  At each iteration, their embedded lengths are evaluated, converted into path-edge weights, and used to form a path-augmented Guttman update.  The process begins from a classical-MDS initialization on the input graph geodesics and repeats until convergence.}
  \label{fig:geodesic_mds_workflow}
\end{figure}

\subsection{Algorithm}

We summarize Geodesic MDS as \cref{alg:geodesic_mds}.

\begin{algorithm}
\caption{Geodesic MDS}
\label{alg:geodesic_mds}
\begin{algorithmic}[1]
  \REQUIRE Data $\bX \in \R^{n \times m}$, target dimension $d$,
           neighborhood size $k$, convergence tolerance $\varepsilon > 0$,
           maximum iterations $T$.
  \ENSURE Embedding $\bZ \in \R^{n \times d}$.
  \\
  \textit{// Step 1: Build input graph geodesic}
  \STATE Construct symmetric $k$-NN graph $\Gk$ on $\bX$ with Euclidean
         edge weights.
  \STATE Compute all-pairs shortest paths $\delta^{(k)}_{ij}$ on $\Gk$
         (Dijkstra / Floyd-Warshall).
  \STATE Choose and store a canonical shortest path family
         $\Gamma_k = \{\gamma_{ij}^{(k)} : i < j\}$ on $\Gk$.
  \\
  \textit{// Step 2: Initialize}
  \STATE Initialize $\bZ^{(0)}$ via a baseline spectral/layout seed
         (e.g.\ classical MDS on $\Delta^{(k)}$ or a multiscale graph layout seed).
  \STATE $t \leftarrow 0$.
  \\
  \textit{// Step 3: Geodesic SMACOF iterations}
  \REPEAT
    \STATE Evaluate the embedded path lengths
           $\demb_{\Gamma_k}(i,j;\bZ^{(t)})$ along the stored geodesics
           $\gamma_{ij}^{(k)}$.
    \STATE Accumulate path-augmented Guttman matrices
           $\tilde{\bV}^{(t)}$, $\tilde{\bW}^{(t)}$ via \cref{eq:path_edge_weight,eq:augmented_guttman}.
    \STATE Update $\bZ^{(t+1)} \leftarrow (\tilde{\bW}^{(t)})^+ \tilde{\bV}^{(t)} \bZ^{(t)}$.
    \STATE $t \leftarrow t + 1$.
  \UNTIL{$\abs{\sigma^{\mathrm{geo}}(\bZ^{(t)}) - \sigma^{\mathrm{geo}}(\bZ^{(t-1)})} < \varepsilon$
         or $t \geq T$.}
  \RETURN $\bZ^{(t)}$.
\end{algorithmic}
\end{algorithm}

\subsubsection*{Initialization}
Classical MDS (eigendecomposition) applied to $\Delta^{(k)}$ is the simplest
baseline initializer and remains a strong candidate on several families.  The
current implementation work reported below shows, however, that it is not a
universal best choice: on highly regular meshes alternative seeds such as KK,
GRIP, Weighted GRIP, or seeded MISF scaffolds can lead to substantially better
final GMDS embeddings.  In practice we therefore recommend treating classical
MDS as one member of an initializer pool rather than as a single default.

\subsubsection*{Multiple restarts}
Since the geodesic stress landscape is non-convex, multiple restarts are
advisable.  The current implementation suggests that the most useful restarts
are often \emph{structured} rather than purely random: spectral seeds, KK
seeds, GRIP-family seeds, and multiscale coarse scaffolds explore genuinely
different basins, whereas small random perturbations of a single seed often do
not.

\subsubsection*{Selection of $k$}
The parameter $k$ controls the tradeoff between tree-metric fidelity
(small $k$) and ambient-metric fidelity (large $k$).  We recommend
selecting $k$ by a cross-validation procedure: hold out a random fraction
of edges from $\Gk$, compute $\delta^{(k)}_{ij}$ on the remaining graph
for held-out pairs, and evaluate how well the embedding reproduces these
held-out geodesics.  The optimal $k$ minimizes held-out geodesic stress.
An alternative is to use the Gromov-hyperbolicity $\delta_0$ of the input
metric as a guide: for nearly tree-like metrics (small $\delta_0$), small
$k$ is appropriate; for nearly Euclidean metrics (large $\delta_0$), large
$k$ is preferable \citep{gromov1987hyperbolic}.

\subsection{A MISF-based Multiscale Modification of GMDS}
\label{sec:misf_gmds}

The experiments of \cref{sec:properties} show that a global classical-MDS
initialization can place GMDS into an unfavorable basin, especially on highly
regular meshes.  A natural alternative is to borrow the
\emph{maximal-independent-set filtration} (MISF) strategy introduced for large
force-directed layouts by \citet{gajer2004multidimensional,gajer2002grip} and
replace the force-directed objective at each scale by the geodesic objective
developed above.  The resulting method is still a fixed-geodesic GMDS method:
the fidelity term remains \cref{eq:raw_geo_stress}, but the initialization and
the subsequent optimization are carried out in a coarse-to-fine manner.

\subsubsection*{Filtration}
Given a connected weighted graph $G=(V,E,w)$, let
\[
  V = V_0 \supset V_1 \supset \cdots \supset V_L
\]
be a maximal-independent-set filtration in the sense of
\citet{gajer2004multidimensional,gajer2002grip}: each $V_{\ell+1}$ is a maximal
subset of $V_\ell$ whose vertices are separated in the graph by a prescribed
minimum graph distance, so that the levels become progressively coarser.
Intuitively, $V_L$ is a sparse sample of the original graph, while $V_0=V$ is
the full vertex set.  The filtration is shown on the underlying $2$D mesh
geometry in \cref{fig:misf_filtration_cases}.  For the regular $12\times 12$
paraboloid the extracted level sizes are $20 \rightarrow 50 \rightarrow 144$;
for the regular $15\times 15$ paraboloid they are
$13 \rightarrow 28 \rightarrow 84 \rightarrow 225$; and for the irregular
rectangle paraboloid they are $11 \rightarrow 32 \rightarrow 81 \rightarrow
225$.

\subsubsection*{Top level}
Instead of initializing the whole graph by classical MDS, we first solve a
\emph{pure} GMDS problem only on the coarsest set $V_L$.  Concretely, we form
the complete graph on $V_L$ with edge weights equal to the graph-geodesic
distances $d_G(i,j)$ measured in the original graph $G$, choose a fixed
geodesic family $\Gamma_L$ on this induced metric graph, and minimize the raw
geodesic stress on $V_L$.  This coarse solve is the \emph{initial embedding} of
the MISF pipeline.  The middle panel in the top row of
\cref{fig:misf_stages_regular_12x12,fig:misf_stages_regular_15x15,fig:misf_stages_irregular_15x15}
shows this initial coarse embedding.  It is not a random seed: it is already
the result of a restartable pure-GMDS solve on $V_L$.  The right panel in the
same row shows that this coarse embedding is later revisited by the sparse
refinement stage.

\subsubsection*{Geodesic insertion of finer levels}
Once $V_{\ell+1}$ has been embedded, each new vertex $v \in V_\ell \setminus
V_{\ell+1}$ is placed with respect to a small anchor set
$A_\ell(v) \subseteq V_{\ell+1}$ drawn from the already embedded coarser level.
The initial location of $v$ is obtained by minimizing
\[
  \Phi_\ell(\bz_v)
  =
  \sum_{a \in A_\ell(v)}
  \omega_{va}
  \bigl( \norm{\bz_v - \bz_a} - d_G(v,a) \bigr)^2,
\]
where the anchor weights $\omega_{va}$ favor well-spread, geodesically close
anchors.  This is a local nonlinear least-squares problem and plays the role of
geodesic trilateration.  In the current implementation each inserted vertex is
optimized independently by a small Armijo line-search solver.  The middle
panels in the lower rows of
\cref{fig:misf_stages_regular_12x12,fig:misf_stages_regular_15x15,fig:misf_stages_irregular_15x15}
show these insertion states before any active-set refinement.

\subsubsection*{Sparse active-set refinement}
After all vertices of level $V_\ell$ have been inserted, the active set is
refined by a sparse GMDS solve.  Let $P_\ell \subseteq V_\ell \times V_\ell$ be
a sparse pair set consisting of short-range geodesic neighbors together with a
small number of landmark pairs.  The level objective has the form
\[
  F_\ell(\bZ)
  =
  \sum_{(i,j)\in P_\ell}
  \bigl(
    \demb_{\Gamma}(i,j;\bZ) - d_G(i,j)
  \bigr)^2
  +
  \lambda_\ell
  \sum_{u \notin V_\ell}
  \alpha_u \norm{\bz_u - \bar{\bz}_u}^2,
\]
where $\bar{\bz}_u$ denotes the current scaffold position of an inactive
vertex.  The anchor term pins the already constructed coarser scaffold while
allowing the active vertices to adjust.  In other words, the method uses the
same fixed-geodesic discrepancy as GMDS, but only on a sparse set of active
pairs, and with the non-active scaffold softly constrained.  The right panels
of \cref{fig:misf_stages_regular_12x12,fig:misf_stages_regular_15x15,fig:misf_stages_irregular_15x15}
show the corresponding refinement states.  After the levelwise cycle reaches
$V_0$, an optional short full-graph pure-GMDS polish is applied.

\subsubsection*{Implementation in \texttt{grip}}
The current prototype reuses the GRIP infrastructure already present in the
\texttt{grip} package.  The MISF itself is extracted by
\texttt{grip.build.misf()} in \texttt{R/grip\_misf.R}, which delegates to the
compiled GRIP filtration code in \texttt{src/DrawGraph.cpp} and
\texttt{src/grip\_rcpp.cpp}.  The multiscale GMDS preparation and orchestration
are implemented in \texttt{R/grip\_geodesic\_misf.R} through
\texttt{grip.prepare.misf.geodesic.mds()} and
\texttt{grip.optimize.misf.geodesic.mds()}.  The per-vertex insertion kernel is
implemented in \texttt{src/geodesic\_misf\_rcpp.cpp}, while both the coarse
top-level solve and the sparse refinement stages reuse the compiled GMDS kernel
in \texttt{src/geodesic\_mds\_rcpp.cpp}.  Thus the MISF modification changes the
optimization schedule and the initialization strategy, but not the underlying
fixed-path geodesic discrepancy.

\subsubsection*{What the multiscale figures show}
\Cref{fig:misf_stages_regular_12x12} shows the regular $12\times 12$
paraboloid.  The coarsest initial embedding on $V_2$ is already far from a
recognizable paraboloid sample ($\rho = 0.8450$), and the subsequent
refinement improves it only to $\rho = 0.5735$.  At the full level $V_0$, the
insertion state is highly tangled and refinement reduces the active-set error
from $\rho = 0.5685$ to $\rho = 0.4543$, but the result is still not a smooth
paraboloid.  \Cref{fig:misf_stages_regular_15x15} shows the same phenomenon on
the regular $15\times 15$ mesh: every level improves under refinement, yet the
coarse sample and the full-level embedding remain visibly far from the target
surface.

By contrast, \cref{fig:misf_stages_irregular_15x15} shows that the irregular
rectangle paraboloid is substantially easier for the same multiscale pipeline.
Its coarsest initial embedding starts much closer to the target
($\rho = 0.2994$), and each subsequent level remains more consistent with the
reference surface.  The final level improves from $\rho = 0.2167$ after
insertion to $\rho = 0.1950$ after sparse refinement, followed by a final full
polish at $\rho = 0.1837$.  This mirrors the empirical observation from the
previous section: the difficulty is not merely the presence of curvature, but
the interaction between the optimization landscape and the high symmetry of the
regular rectangular mesh.

\begin{table}[t]
  \centering
  \small
  \caption{Stagewise active-set Procrustes error $\rho$ for the MISF-GMDS
  pipeline on the three paraboloid families shown in
  \cref{fig:misf_stages_regular_12x12,fig:misf_stages_regular_15x15,fig:misf_stages_irregular_15x15}.  The
  ``initial'' column is the coarse top-level GMDS solve for the first active
  level and the anchor-based insertion state for the remaining levels; the
  ``refined'' column is the corresponding sparse GMDS refinement.}
  \label{tab:misf_stage_rho}
  \begin{tabular}{llrrr}
    \toprule
    case & level & $|V_\ell|$ & initial $\rho$ & refined $\rho$ \\
    \midrule
    regular $12\times 12$ & $V_2$ & 20  & 0.8450 & 0.5735 \\
    regular $12\times 12$ & $V_1$ & 50  & 0.6810 & 0.5431 \\
    regular $12\times 12$ & $V_0$ & 144 & 0.5685 & 0.4543 \\
    regular $12\times 12$ & final & 144 & n/a    & 0.4546 \\
    \midrule
    regular $15\times 15$ & $V_3$ & 13  & 0.9268 & 0.7982 \\
    regular $15\times 15$ & $V_2$ & 28  & 0.8841 & 0.7366 \\
    regular $15\times 15$ & $V_1$ & 84  & 0.7193 & 0.6106 \\
    regular $15\times 15$ & $V_0$ & 225 & 0.6870 & 0.6308 \\
    regular $15\times 15$ & final & 225 & n/a    & 0.6174 \\
    \midrule
    irregular $15\times 15$ & $V_3$ & 11  & 0.2994 & 0.2281 \\
    irregular $15\times 15$ & $V_2$ & 32  & 0.2741 & 0.2349 \\
    irregular $15\times 15$ & $V_1$ & 81  & 0.2345 & 0.1987 \\
    irregular $15\times 15$ & $V_0$ & 225 & 0.2167 & 0.1950 \\
    irregular $15\times 15$ & final & 225 & n/a    & 0.1837 \\
    \bottomrule
  \end{tabular}
\end{table}

\begin{figure}[t]
  \centering
  \includegraphics[width=\linewidth]{../figures/fig_misf_filtration_cases.pdf}
  \caption{MIS filtration levels displayed on the underlying $2$D mesh geometry.  Each
  row shows one graph family, with the active level $V_\ell$ highlighted on top
  of the full mesh graph.  The irregular rectangle keeps rectangular topology
  but breaks the lattice symmetry in the geometry, yielding a more generic
  coarse sample than the perfectly regular meshes.}
  \label{fig:misf_filtration_cases}
\end{figure}

\begin{figure}[t]
  \centering
  \includegraphics[width=\linewidth]{../figures/fig_misf_stages_regular_12x12.pdf}
  \caption{Regular $12\times 12$ paraboloid under MISF-GMDS.  Each row
  corresponds to one active set $V_\ell$.  Left: the reference active-set
  sample.  Middle: the initial multiscale embedding for that level (the coarse
  pure-GMDS solve for $V_2$, then anchor-based insertion for $V_1$ and $V_0$).
  Right: the sparse GMDS refinement of the same active set.  The coarsest level
  is already not paraboloid-like, and although every level improves, the full
  embedding remains folded.}
  \label{fig:misf_stages_regular_12x12}
\end{figure}

\begin{figure}[t]
  \centering
  \includegraphics[width=\linewidth]{../figures/fig_misf_stages_regular_15x15.pdf}
  \caption{Regular $15\times 15$ paraboloid under MISF-GMDS.  The same
  qualitative pattern persists at the larger size: the coarsest-level solve does
  not recover a recognizable paraboloid sample, insertion creates a tangled
  full-level state, and sparse refinement improves but does not remove the
  pathology.}
  \label{fig:misf_stages_regular_15x15}
\end{figure}

\begin{figure}[t]
  \centering
  \includegraphics[width=\linewidth]{../figures/fig_misf_stages_irregular_rectangle_15x15.pdf}
  \caption{Irregular $15\times 15$ rectangle lifted to a paraboloid.  The
  multiscale pipeline is much better behaved on this less symmetric graph: the
  coarsest level already resembles a sparse paraboloid sample, and each finer
  refinement remains visibly closer to the reference geometry than in the
  regular-grid cases.}
  \label{fig:misf_stages_irregular_15x15}
\end{figure}

\subsection{Relationship to Metric Geometry}
\label{sec:metric_geometry}

The construction of Geodesic MDS is grounded in several concepts from
metric geometry that illuminate its properties.

\subsubsection*{Minimum realization}
The $k$-NN graph $\Gk$, after removing metrically redundant edges
(pairs $(i,j)$ for which $\delta^{(k)}_{ij}$ equals the shortest path
through other vertices), yields the \emph{minimum realization} of the
input graph metric \citep{dress1984trees, imrich1984optimal}.
Geodesic MDS with this minimum realization as the underlying graph
uses the most parsimonious geodesic structure consistent with the data.

\subsubsection*{Tight span}
The \emph{tight span} $T(\Delta)$ (or injective hull)
\citep{dress1984trees, isbell1964six} is the minimal injective metric
space containing $(X, \delta^{(k)})$.  It is a polyhedral complex whose
1-skeleton contains the minimum realization.  Geodesic MDS seeks an
embedding of the 0-skeleton (the $n$ data points) into $\R^d$ that
preserves the 1-skeleton path metric, which is a projection of the
tight span onto the original vertex set.

\subsubsection*{Metric spanners}
As noted in the background discussion of the $k$-NN graph, $\Gk$ is a $(1+\varepsilon)$-spanner
of the input Euclidean metric.  Geodesic MDS with a spanner graph thus
approximates the full Euclidean MDS problem with a multiplicative error
of $(1+\varepsilon)$ on input dissimilarities, while using a sparse graph
structure that enforces geodesic consistency \citep{althofer1993sparse}.

\subsubsection*{Gromov-Hausdorff approximation}
The graph metric $(\{1,\ldots,n\}, d_{\Gk})$ converges in the
Gromov-Hausdorff sense to the underlying geodesic metric space as
$n \to \infty$ and $k$ grows appropriately
\citep{gromov1999metric, latschev2001riemannian, bernstein2000graph}.
Geodesic MDS can therefore be interpreted as finding an optimal isometric
embedding of a Gromov-Hausdorff approximation of the data geometry, rather
than an optimal embedding of the ambient Euclidean point cloud.

\subsection{Current State of the Working Implementation}
\label{sec:current_implementation}

The working implementation in the \texttt{grip} package is now best viewed as
a family of related GMDS prototypes rather than a single finished algorithm.
The common core is still the fixed-geodesic discrepancy
\cref{eq:raw_geo_stress}: for a prepared graph-geodesic family
$(G,\Gamma)$, the implementation measures the embedded length of the
input-side geodesics and optimizes that mismatch in a compiled C++ kernel.
Around that core, several experimental layers have been added in order to
understand the main failure modes of the objective.

\begin{table}[t]
  \centering
  \small
  \caption{Current GMDS implementation layers in \texttt{grip}.}
  \label{tab:implementation_layers}
  \begin{tabular}{p{0.23\linewidth} p{0.34\linewidth} p{0.35\linewidth}}
    \toprule
    layer & main files / entry points & current role \\
    \midrule
    Pure and regularized GMDS &
    \texttt{R/grip\_quality.R},
    \texttt{src/geodesic\_mds\_rcpp.cpp},
    \texttt{grip.prepare.graph.geodesic.mds()},
    \texttt{grip.optimize.geodesic.mds()} &
    Fixed-geodesic objective, tie-averaged shortest paths, anchor regularization, graph-distance-aware repulsion, and edge-spring relaxations. \\
    \addlinespace
    MISF-based GMDS &
    \texttt{R/grip\_misf.R},
    \texttt{R/grip\_geodesic\_misf.R},
    \texttt{src/geodesic\_misf\_rcpp.cpp},
    \texttt{grip.optimize.misf.geodesic.mds()} &
    Coarse-to-fine GMDS pipeline with top-level solve, per-vertex insertion, sparse level refinement, and final full-graph polish. \\
    \addlinespace
    Layout-based basin finders &
    \texttt{R/grip\_layout\_weighted.R},
    \texttt{src/MishWeighted.cpp},
    \texttt{src/grip\_rcpp.cpp},
    Weighted GRIP / weighted KK / LGKK APIs &
    External basin-finding and warm-start machinery for MISF-GMDS; currently the most useful source of coarse scaffolds and lower-level placement states. \\
    \bottomrule
  \end{tabular}
\end{table}

At this stage the implementation already supports several experimentally
distinct variants: pure GMDS; tie-averaged GMDS; anchor-regularized GMDS;
anchor-plus-repulsion and anchor-plus-edge-spring regularizations; an
edge-relaxation surrogate; a MISF-based multiscale GMDS schedule; and seeded
MISF variants that reuse KK, GRIP, and Weighted GRIP layouts as basin
finders.  The goal of the recent experimental cycle was not merely to optimize
one more hyperparameter, but to determine \emph{which part of the current
pipeline is actually responsible for the observed pathologies}.

\subsection{What the Recent Experiments Show}
\label{sec:recent_experiments}

\subsubsection*{The smooth target may exist even when the optimizer misses it}
The identifiability experiment on the regular and irregular paraboloid meshes
showed that the smooth paraboloid is an \emph{exact} GMDS solution after the
graph-weight normalization used by the generator.  In particular, for the
regular $12\times 12$ paraboloid, a $15^\circ$ folded perturbation of the
truth has only moderate raw geodesic stress
($\sigma = 0.0102$, $\rho = 0.0325$), and pure GMDS started from that folded
state returns close to the smooth reference
($\sigma = 0.0002$, $\rho = 0.0303$ after $25$ iterations).  By contrast,
pure GMDS started from a global classical-MDS seed lands in a very different
basin ($\sigma = 0.0298$, $\rho = 0.3291$ on the same graph).  Thus the main
issue is not absence of a smooth solution but basin selection.

\subsubsection*{Single-scale regularization helps locally but does not solve the
basin problem}
On the regular paraboloid meshes, anchor-based regularization improves the raw
geodesic objective substantially but leaves the recovered shape much too far
from the target.  For the regular $12\times 12$ paraboloid, direct
\texttt{cmdscale $\rightarrow$ pure GMDS} ends at
$(\sigma,\rho)=(0.0492, 0.3257)$, while the best single-scale regularized run
reaches $(0.0221, 0.3214)$.  For the regular $15\times 15$ paraboloid, the
best regularized run reaches $(0.0183, 0.2286)$, improving stress much more
than shape.  In other words, the current anchor, repulsion, and edge-spring
terms are useful as controlled perturbations of the objective, but they do not
yet provide a reliable route to the correct basin on their own.

\subsubsection*{Multiscale GMDS helps only if the scaffold is right}
The first MISF-GMDS prototype in \cref{sec:misf_gmds} showed that a pure-GMDS
coarse solve is often not a good top-level scaffold on regular meshes.  The
subsequent Phase~A--E experiments therefore replaced ``one GMDS initializer''
by a seed pool and asked which part of the multiscale schedule actually needs
help.  The answer is now fairly clear.  The top MISF level is a
basin-selection problem, not merely a local refinement problem; lower-level
geodesic insertion can destroy a good coarse scaffold; weighted KK is
currently the strongest lower-level placement mechanism; and a small pool of
top-level seeds is better than any single universal default.

\begin{table}[t]
  \centering
  \small
  \caption{Representative recent outcomes on the paraboloid families.  The
  ``direct'' column is the full-graph \texttt{cmdscale $\rightarrow$ pure GMDS}
  control.  The ``best single-scale regularized'' column reports the best
  regularized-GMDS run from the single-scale study.  The last column gives the
  best currently observed seeded multiscale pipeline on the corresponding
  family.}
  \label{tab:recent_paraboloid_results}
  \begin{tabular}{lccc}
    \toprule
    case & direct $(\sigma,\rho)$ & best single-scale regularized $(\sigma,\rho)$ & current best multiscale $(\sigma,\rho)$ \\
    \midrule
    regular paraboloid $12\times 12$ & $(0.0492, 0.3257)$ & $(0.0221, 0.3214)$ & $(0.0109, 0.0904)$ \\
    regular paraboloid $15\times 15$ & $(0.0519, 0.2317)$ & $(0.0183, 0.2286)$ & $(0.0148, 0.0726)$ \\
    irregular rectangle paraboloid $15\times 15$ & $(0.2872, 0.1845)$ & n/a & $(0.0951, 0.0889)$ \\
    \bottomrule
  \end{tabular}
\end{table}

\begin{figure}[t]
  \centering
  \includegraphics[width=0.98\linewidth]{../figures/fig_current_implementation_summary.pdf}
  \caption{Summary of the current implementation state.  Left: progression of
  final Procrustes error $\rho$ on the three paraboloid families as the
  pipeline moves from direct full-graph GMDS to seeded multiscale variants.
  Right: Block~1 of the Phase~E mesh-geometry study shows that the current
  truth-free selector is still too \texttt{cmdscale}-biased: it chose
  \texttt{cmdscale} in $10/12$ cases, whereas the oracle best final-$\rho$
  seed was \texttt{kk} in $6/12$ cases, \texttt{cmdscale} in only $3/12$,
  \texttt{grip} in $2/12$, and \texttt{weighted\_grip\_polish\_lgkk} in
  $1/12$.}
  \label{fig:current_implementation_summary}
\end{figure}

\subsubsection*{The current selector is not yet good enough}
The most recent Block~1 mesh-geometry expansion confirms that there is no
single universal top-level winner across flat, saddle, paraboloid, and ripple
meshes.  The current truth-free proxy chose \texttt{cmdscale} in $10/12$
cases and matched the oracle best final-$\rho$ seed in only $4/12$ cases.
This is precisely why the current pipeline should still be viewed as a
research prototype: the best-performing top seed depends on geometry, and the
automatic selector still over-rewards immediate low-stress spectral states.

\subsection{Immediate Next Steps}
\label{sec:next_steps}

The recent experiments reduce the open problem to a much narrower one than we
had at the beginning of the project.  The main challenge is no longer to
decide whether GMDS should be single-scale or multiscale, or whether classical
MDS should be discarded entirely.  Instead, the present evidence suggests the
following agenda.

\subsubsection*{1. Improve the truth-free seed selector}
The current Phase~E proxy uses only a one-level weighted-KK lookahead together
with three rank terms (lookahead stress, roughness, and edge-floor).  This is
not yet predictive enough.  The next selector experiments should include
insertion-stability signals such as area-floor diagnostics, first-level energy
jumps, and possibly a second short refinement step before choosing the winner.
An especially promising direction is to retain a small shortlist of seeds
rather than forcing a single winner too early.

\subsubsection*{2. Keep weighted KK as the current lower-level default}
Across the recent multiscale studies, weighted KK is the strongest lower-level
placement rule we currently have.  It should therefore remain the default
lower-level mechanism while the top-level seed-selection problem is being
improved.  Future comparisons should treat GRIP-family lower-level placement as
an ablation against this weighted-KK baseline, not as the default.

\subsubsection*{3. Expand only after the selector improves}
The ongoing mesh-geometry expansion is useful because it reveals where the
selector fails, but the next full cross-family expansion should wait until the
top-level proxy is better calibrated.  Otherwise we risk spending substantial
computation on a pipeline whose automatic choice rule is not yet trustworthy.
The most natural order is: first improve the selector on the mesh geometries,
then move to perforated rectangles, and only then to annulus/sphere families.

\subsubsection*{4. Continue objective-level diagnostics without conflating them with
the optimizer}
The experiments already show that the smooth target may be an exact GMDS
solution while the current initialization misses it.  This means that
objective-level questions (fold-blindness, path averaging, repulsion,
edge-relaxation) and optimization-level questions (basin finding, multiscale
placement, seed selection) should continue to be studied in parallel but kept
conceptually separate.  At present, the most urgent unresolved issue is
optimization/basin selection rather than the formal definition of GMDS itself.

%% ============================================================
%% BIBLIOGRAPHY
%% ============================================================

\section*{Acknowledgments}
Research reported in this publication was supported by grant number
INV-048956 from the Gates Foundation.

\bibliographystyle{plainnat}
\bibliography{geodesic_mds}

\end{document}
